\documentclass{template}

\usepackage{mathtools}
\pgfplotsset{compat=1.18}

\newcommand{\qedwhite}{\hfill \ensuremath{\Box}}


\begin{document}

\begin{center}
    \Large \textcolor{blue!70!black}{\textbf{Commutative Algebra Problem Set \#1}} \\
    \normalsize \textit{Dmytro Zakharov}
\end{center}

\begin{problem}[Exercise 1]
    Let $\mathfrak{p}$ and $\mathfrak{q}$ be two prime ideals in a ring $A$ such
    that neither $\mathfrak{p} \subseteq \mathfrak{q}$ nor $\mathfrak{q}
    \subseteq \mathfrak{p}$. Show that the ideal $\mathfrak{p} \cap
    \mathfrak{q}$ is not prime.
\end{problem}

\textbf{Solution.} To show that $\mathfrak{p} \cap \mathfrak{q}$ is not prime,
it suffices to pick two $a,b$ such that $ab \in \mathfrak{p} \cap \mathfrak{q}$,
but $a,b \not\in \mathfrak{p} \cap \mathfrak{q}$. Since neither $\mathfrak{p}
\subseteq \mathfrak{q}$ nor $\mathfrak{q} \subseteq \mathfrak{p}$ holds, there
exist elements $a \in \mathfrak{p} \setminus \mathfrak{q}$ and $b \in
\mathfrak{q} \setminus \mathfrak{p}$. Notice that both elements are not in
$\mathfrak{p} \cap \mathfrak{q}$, but their product is: on one hand, $ab \in
\mathfrak{p}$ since $a \in \mathfrak{p}$ and $\mathfrak{p}$ is an ideal.
Similarly, $ab \in \mathfrak{q}$ since $b \in \mathfrak{q}$. Thus, since $ab \in
\mathfrak{p}$ and $ab \in \mathfrak{q}$, we conclude $ab \in \mathfrak{p} \cap
\mathfrak{q}$. \qedwhite

\begin{problem}[Exercise 2]
    A ring $A$ is Boolean if $x^2=x$ for all $x \in A$. In a 
    Boolean ring $A$, show that:
    \begin{itemize}
        \item[\textit{(i)}] $2x=0$ for all $x \in A$;
        \item[\textit{(ii)}] every prime ideal $\mathfrak{p}$ is 
        maximal, and $A/\mathfrak{p}$ is a field with two elements;
        \item[\textit{(iii)}] every finitely generated ideal in 
        $A$ is principal. 
    \end{itemize}
\end{problem}

\textbf{Solution.} 

\textit{Part (i).} For arbitrary $x \in A$, we conduct the following steps:
\begin{equation*}
    2x = 4x - 2x = 4x^2 - 2x = (2x)^2 - 2x = 2x - 2x = 0
\end{equation*}

\textit{Part (ii).} Take arbitrary $[x] \in A/\mathfrak{p}$ (we denote by
$[\cdot]$ elements of quotient ring), where $\mathfrak{p}$ is a prime ideal.
Note that $[x]^2=[x^2]=[x]$, thus we have $[x]^2-[x] = ([x]-[1])[x] = [0]$.
Since $\mathfrak{p}$ is prime, $A/\mathfrak{p}$ is a domain, so either $[x]=[1]$
or $[x]=[0]$. This proves that $A/\mathfrak{p} = \{[0],[1]\}$ is a domain with
two elements. However, it is clearly a field: the only non-zero element is $[1]$
and its inverse is simply $[1]$. Therefore, $A/\mathfrak{p} \cong \mathbb{F}_2$.

\textit{Part (iii).} It suffices to prove this fact for an ideal generated by
two elements: say, $\mathfrak{m} = \langle a_1, a_2\rangle$. Namely, we prove
that $\langle a_1, a_2 \rangle = \langle a \rangle$ for $a=a_1+a_2-a_1a_2$. The
general case will follow by applying induction. 

\textcolor{gray!80!black}{$a \in \langle a_1, a_2 \rangle$.} This inclusion is
evident: $a$ is a linear combination of $a_1$ and $a_2$.

\textcolor{gray!80!black}{$a_1,a_2 \in \langle a \rangle$.} Note that
\begin{equation*}
    aa_1 = (a_1+a_2-a_1a_2)a_1 = a_1^2 + a_2a_1 - a_1^2a_2 = a_1
\end{equation*}

Thus $a_1 \in \langle a \rangle$. Similarly, $a_2 \in \langle a \rangle$. Thus
we have shown $\langle a_1, a_2 \rangle = \langle a \rangle$. \qedwhite

\begin{problem}[Exercise 3]
    This exercises consists of two subproblems:
    \begin{itemize}
        \item[\textit{(i)}] Show that a ring $A$ is a domain if and only if it
admits an injective homomorphism $A \hookrightarrow K$ into a field $K$.
        \item[\textit{(ii)}] Show that $A$ is reduced (has no non-zero nilpotent
elements) if and only if it admits an injective homomorphism $A \hookrightarrow
\prod_{\alpha \in I}K_{\alpha}$ into a product of (possibly infinitely many)
fields.
    \end{itemize}
\end{problem}

\textbf{Solution.}

\textit{Part (i).} $(\Rightarrow)$ Suppose $A$ is a domain. Now consider the
natural inclusion morphism $\iota: A \hookrightarrow \text{Frac}(A), a \mapsto a/1$. Clearly,
this is an injective homomorphism, and $K=\text{Frac}(A)$ is indeed a field.

$(\Leftarrow)$ Suppose $K$ is a field, and one is provided with the injective
homomorphism $\gamma: A \hookrightarrow K$. We shall show that $A$ is a domain. Take two
elements $x,y \in A$ such that $xy=0$. Take homomorphism of both sides:
$\gamma(x)\gamma(y)=\gamma(0)$. Clearly, $\gamma(0)=0$, so we are left with
$\gamma(x)\gamma(y) = 0$. Since $K$ is a field (and an integral domain, in 
particular), either $\gamma(x) = 0$ or $\gamma(y) = 0$. But since $\gamma$
is injective, this implies that either $x=0$ or $y=0$. Thus $A$ is 
a domain. \qedwhite

\textit{Part (ii).} $(\Rightarrow)$  Recall that $\text{Nil}(A) =
\bigcap_{\mathfrak{p} \; \text{prime}}\mathfrak{p}$. Since $A$ is reduced,
$\text{Nil}(A) = 0$, so no two prime ideals intersect besides $0$. This gives
the natural embedding $A \hookrightarrow \prod_{\alpha \in I}A/\mathfrak{p}$ by
sending $x \in A$ to the corresponding quotient $A/\mathfrak{p}$. But note that
since $\mathfrak{p}$ is prime, $A/\mathfrak{p}$ is a domain, and thus we have a
natural embedding into $\text{Frac}(A/\mathfrak{p})$. Composing both embeddings 
produce the desired result.

$(\Leftarrow)$ Suppose we are given injective homomorphism $\varphi: A
\hookrightarrow \prod_{\alpha \in I}K_{\alpha}$. Let $a \in A$ be nilpotent, so
$a^n = 0$ for some $n$. Apply $\varphi$ to obtain $\varphi(a^n) =
\varphi(a)^n=0$. This means that $\varphi(a)_{\alpha}^n=0$ for each $\alpha \in
I$. Since field has no non-zero nilpotents, we conclude that
$\varphi(a)_{\alpha}=0$ or simply $\varphi(a)=0$. Since $\varphi$ is injective,
$a=0$. Thus, there are no non-zero nilpotents in $A$, so $A$ is reduced. \qedwhite

\begin{problem}[Exercise 4]
    Let $X$ be a compact Hausdorff space and let $A=C(X,\mathbb{R})$ be the ring
of continuous functions on $X$. For $x \in X$, let $\mathfrak{m}_x \subseteq A$
be the set of all $f \in A$ such that $f(x)=0$. Show that $\mathfrak{m}_x$ is a
maximal ideal in $A$, and that every maximal ideal $\mathfrak{m} \subseteq A$ is
of the form $\mathfrak{m}_x$ for a unique $x \in X$.
\end{problem}

\textbf{Solution.} To prove $\mathfrak{m}_x$ is a maximal ideal, we prove that
$A/\mathfrak{m}_x$ is a field. I see two approaches how to prove this. 

\textit{Approach \#1 (mechanical).} Take any $[f] \in A/\mathfrak{m}_x$. If $f
\in \mathfrak{m}_x$ (that is, $f(x)=0$), then $[f] = [0]$. Otherwise, the
inverse of $[f]$ is always given by the constant function representative
$[1/f(x)]$. This function exists since $f(x) \neq 0$, and moreover $[f][1/f(x)]
= [f/f(x)] = [1]$.

\textit{Approach \#2.} Consider the evaluation morphism $\text{ev}_x: A \to
\mathbb{R}$ that sends $f \mapsto f(x)$. Note that $\text{ker}(\text{ev}_x) =
\mathfrak{m}_x$ by definition of the ideal $\mathfrak{m}_x$. Moreover, clearly
$\text{im}(\text{ev}_x) = \mathbb{R}$ (since for any $\alpha \in \mathbb{R}$ we
can find some continuous function $f \in A$ such that $f(x)=\alpha$). Thus, we
induce the isomorphism $A/\mathfrak{m}_x \cong \mathbb{R}$. Since $\mathbb{R}$
is a field, $A/\mathfrak{m}_x$ is a field. \qedwhite

Now we prove that all maximal ideals are of form $\mathfrak{m}_x$ with $x \in
X$. For this, we will show the auxiliary proposition.

\begin{propositions}
    Let $\mathfrak{m}$ be a maximal ideal and $V(\mathfrak{m}) := \{x \in X:
f(x) = 0 \; \text{for all} \; f \in \mathfrak{m}\}$. Then, $V(\mathfrak{m}) =
\{x\}$ for some $x \in X$, and $\mathfrak{m} \subseteq \mathfrak{m}_x$.
\end{propositions}

\textbf{Proof.} We first show that this set is non-empty. Indeed, suppose
otherwise: $V(\mathfrak{m}) = \emptyset$, thus for each $x \in X$ we can find
$f_x \in \mathfrak{m}$ such that $f_x(x) \neq 0$. Thus, by continuity, these
exists the neighbourhood $U_x$ around $x$ such that $f_x(U_x) \neq \{0\}$. Note
that $\bigcup_{x \in X}U_x$ covers the whole $X$. Since $X$ is a compact space,
we might choose the finite subcovering: for instance, let $\bigcup_{i=1}^n
U_{x_i} = X$. Let $h := \sum_{i=1}^n f_{x_i}^2$. Clearly, $h \in \mathfrak{m}$
since every $f_{x_i} \in \mathfrak{m}$. However, note that $h$ is a strictly
positive function over $X$ since it is a sum of squares of strictly positive
functions over corresponding open sets $U_{x_i}$. Therefore, $h \in
\mathfrak{m}$ is a unit. However, the maximal ideal cannot contain a unit
(otherwise $\mathfrak{m}$ is the whole ring). Contradiction.

Now suppose $x \in V(\mathfrak{m})$. Notice that $\mathfrak{m} \subseteq
\mathfrak{m}_x$. Indeed, by definition of $V(\mathfrak{m})$, every element of $f
\in \mathfrak{m}$ satisfies $f(x) = 0$, so we have the corresponding inclusion.
But since $\mathfrak{m}$ is maximal, either $\mathfrak{m}_x = \mathfrak{m}$ or
$\mathfrak{m}_x = A$. The latter is clearly not the case, so we conclude
$\mathfrak{m}_x = \mathfrak{m}$.

Finally, we show uniqueness. Suppose $x,y \in V(\mathfrak{m})$, so in particular
we have $\mathfrak{m}_x = \mathfrak{m}_y = \mathfrak{m}$. Now we apply the
Tietze's extension theorem: define function $g: \{x,y\} \to \mathbb{R}$ such
that $g(x) = 0, g(y) = 1$. Then, there exists the continuous extension
$\widetilde{g} \in A$ such that $\widetilde{g}(x)=0$ and $\widetilde{g}(y)=1$.
Since $\widetilde{g}(x)=0$, $\widetilde{g} \in \mathfrak{m}_x$, but that implies
$\widetilde{g} \in \mathfrak{m}_y$. But since $\widetilde{g}(y)=1$, this is a
contradiction.

Notice that this proposition finishes the theorem off. \qedwhite

\begin{problem}[Exercise 5]
    Let again $X$ be a compact Hausdorff space and let $A=C(X,\mathbb{R})$ be
the ring of continuous functions on $X$. Show that every prime ideal in $A$ is
contained in a unique maximal ideal.
\end{problem}

\textbf{Solution.} We apply the same idea as in the previous proof. For prime
ideal $\mathfrak{p}$, define the vanishing set $V(\mathfrak{p}) = \{x \in X:
f(x)=0 \; \text{for all} \; f \in \mathfrak{p}\}$. First note that this set is
non-empty for the same very reasons, as in the proposition above (unit cannot be
contained in prime ideal either). Thus, there is some $x \in V(\mathfrak{p})$
and $\mathfrak{p} \subseteq \mathfrak{m}_x$ by the definition of
$V(\mathfrak{p})$. It is remained to show that $x$ is unique. To show that, we
will need the following proposition.
\begin{propositions}
    For $E \subseteq X$, define $I(E) := \{f \in A: f(x) = 0 \; \text{for all}
    \; x \in E\}$. Then, $I(V(\mathfrak{a})) = \sqrt{\mathfrak{a}}$ for any
    ideal $\mathfrak{a} \subseteq A$.
\end{propositions}

For now, assume that this proposition holds (we will prove it below shortly).
Then, assume $x,y \in V(\mathfrak{p})$ and $x \neq y$. Then, by Tietze’s
extension theorem, we can find two functions $f,g \in A$ such that $f(x)=0$,
$f(y)=1$, and $g(x)=1$, $g(y)=0$. Moreover, it is easy to check that we can
construct $f$ and $g$ in such a way that $f(z)g(z)=0$ for every $z \in X$. In
particular, $fg \in I(V(\mathfrak{p}))$, thus $fg \in
\sqrt{\mathfrak{p}}=\mathfrak{p}$. Thus, either $f \in \mathfrak{p}$ or $g \in
\mathfrak{p}$. But $f$ and $g$ clearly do not vanish on $V(\mathfrak{p})$, so we
get the contradiction. Thus, $V(\mathfrak{p})=\{x\}$ for some $x \in X$, which
proves the theorem. \qedwhite

\textbf{Proposition proof.} $\sqrt{\mathfrak{a}} \subseteq I(V(\mathfrak{a}))$
is obvious: take $f \in A$ and we know that $f^n \in \mathfrak{a}$ for some $n$.
For any $x \in V(\mathfrak{a})$ we know that $f^n(x)=0$ by definition of
$V(\mathfrak{a})$. Clearly, this implies $f(x)=0$ and thus $f \in
I(V(\mathfrak{a}))$.

The hard part is showing $I(V(\mathfrak{a})) \subseteq \sqrt{\mathfrak{a}}$. In
other words, suppose $f$ vanishes over all $x \in V(\mathfrak{a})$, then we
shall show that $f^n \in \mathfrak{a}$ for some $n$. Define the set $Y := X
\setminus V(\mathfrak{a})$. This set is open, because
$V(\mathfrak{a})=\bigcap_{f \in \mathfrak{a}} \{x \in X: f(x)=0\}$ is the
intersection of closed sets. Thus, for every $y \in Y$, find some function $f_y
\in \mathfrak{a}$ and neighbourhood $U_y$ such that $f_y(U_y) \neq \{0\}$.
Clearly, $\{U_y\}_{y \in Y}$ is an open covering of $Y$, which means $\{U_y\}_{y
\in Y} \cup V(\mathfrak{a})$ covers all $X$. This means that there exists a
finite covering of $Y$: let it be $\{U_{y_i}\}_{i=1}^n$ with corresponding
functions $f_{y_i}$. Now define $s := \sum_{i=1}^n f_{y_i}^2 \in \mathfrak{a}$.
Note that we constructed $s$ such that:
\begin{equation*}
    s|_{X\setminus V(\mathfrak{a})} > 0, \quad s|_{V(\mathfrak{a})} = 0, \quad s \in \mathfrak{a}.
\end{equation*}

Now define the following function:
\begin{equation*}
    \varphi(x) := \begin{cases}
        f(x)^2/s(x), & \text{if} \; s(x) \neq 0, \\
        0, & \text{if} \; s(x) = 0.
    \end{cases}
\end{equation*}

Notice that $\varphi \in A$ and $f^2 = \varphi s$. Since $s \in \mathfrak{a}$,
we conclude that $\varphi s \in \mathfrak{a}$ and thus $f^2 \in \mathfrak{a}$,
so in particular $f \in \sqrt{\mathfrak{a}}$. \qedwhite

\end{document}